\subsection{全量指标及其权重}
场景 B 五核的 $100$ 图总周期为 $65\,769\,876$，相对固定单核 $312\,479\,373$ 周期下降 $78.9523\%$；总周期比为 $4.751102$。但赛题规定的五核逐图平均加速比为 $4.159974$，两者相差明显。这并不表示某一套结果错算，而是式\eqref{eq:weighted-speedup}所示的加权方式不同：总周期比给仍耗时很长的图更高权重，逐图平均让每张图的速度贡献相同。为避免在论述“效果”时偷换口径，表\ref{tab:q2-metric-perspective}同时保留两列，并只以逐图平均加速比作为规定折线图的点。
\begin{table}[!htbp]\centering
\caption{问题二的逐图平均与百图总周期口径}\label{tab:q2-metric-perspective}
\begin{tabular}{rrrrr}
\toprule
核数 & 平均加速比 & 总周期比 & 总周期下降率 & 总额外搬运（bytes）\\
\midrule
1 & 1.000000 & 1.000000 & $0.0000\%$ & 1\,241\,686\,120\\
2 & 1.981113 & 2.030721 & $50.7564\%$ & 1\,128\,799\,934\\
3 & 2.815714 & 3.005163 & $66.7239\%$ & 1\,079\,358\,470\\
4 & 3.550822 & 3.918488 & $74.4800\%$ & 1\,092\,834\,382\\
5 & 4.159974 & 4.751102 & $78.9523\%$ & 1\,087\,391\,086\\
\bottomrule
\end{tabular}
\end{table}
问题一与问题二都使用同一张原始图和同一固定单核参考，但两问的 Task、跨图搬运与同步规则不同，且各自重新选了方案。因此例如“五核 B 总周期小于 A”是\emph{场景与算法共同变化}后的结果，不能直接宣称“500周期同步比1000周期快多少”或“同核驻留单独贡献多少”。如果要单独归因某个机制，应在同一方案上做合法的成对反事实评价，并保持其余规则与硬件参数相同；本文的问题三正采用了同方案硬件对照来避免这一类混淆。

\subsection{非单调图与特定结构的解释}
增加核心数并不自动扩展同一方案的无损可行域：现行算法对每个 $k$ 分别划分、映射与排序，有界预算产生的候选池也不同。全 $100$ 图、相邻核心数的 $400$ 个比较中，有 $12$ 处较多核配置反而更慢。第64图四核为 $10\,160$ 周期、五核为 $11\,391$ 周期；第71图三核为 $9\,867$ 周期、四核为 $10\,732$ 周期；第44图三核 $44\,484$、四核 $49\,989$ 周期，但五核又降至 $43\,028$。这些点说明“平均曲线单调”与“每图曲线单调”是两个不同命题，也说明重新划分时可能失去先前核数下的有利结构。仅凭周期与搬运汇总尚不足以把每个反例唯一归因为同步、DDR 或缓存；需要固定方案与时间线才能分析具体控制路径。

第5图从单核 $95\,618$ 周期逐渐降至五核 $50\,673$ 周期，同时额外搬运从单核 $0$ 增至五核 $1\,688\,564$ 字节。它体现了主目标与次指标的取舍：如果为了控制搬运而保留整图在一个核心，可能错过可用并行；如果只追求更多任务并行，又会扩大跨核 COPY 和同步等待。当前目标先比较周期，在相同周期下选较少搬运；论文不把这一个图的结果推广为所有图的规律。

\subsection{自身版本对照与结果可信范围}
当前完整 B 方案与上一轮本工程的百图 B 方案在同一 $400$ 个多核配置上比较，$13$ 组周期改善、$387$ 组持平、无退步。四种多核配置合计从 $403\,398\,480$ 降至 $403\,371\,635$ 周期，减少 $26\,845$ 周期，即约 $0.006655\%$；同时逻辑额外搬运合计从 $4\,387\,116\,718$ 增至 $4\,388\,383\,872$ 字节，增加 $1\,267\,154$ 字节，即约 $0.028884\%$。数字说明当前机制有小幅\emph{已观测}收益，也说明收益不是“毫无代价”的。新结构准入、重复机会补位与普通联合邻域共同组成这一次版本变化，现有全量对照不支持把 $26\,845$ 周期分摊给单一模块。

对全部 $500$ 个配置，保存的方案均可在官方 B 规则下重放；$400$ 个多核配置的物理张量容量审计与完整结果字段核对通过。单核 $100$ 个点采用题目规定的整图核内启发式参考，而不是从多核最优候选回推出单核。这样的复核证明结果与本次数据、配置和官方评估语义一致，却不是全局最优证书或对未见计算图的统计泛化检验。若更换核数上限、DDR 带宽或内存容量，结构门控、局部成本和候选排序均须重新检查。
